
Uncertainty-based hallucination detection methods report strong results in isolation---semantic entropy achieves AUROC values in the range 0.75--0.85 on TruthfulQA according to Kuhn et al. (2023), and self-consistency reaches approximately 0.70--0.80 on WikiBio according to Manakul et al. (2023). However, these methods have been evaluated on different benchmarks with different sample counts and different models. This pilot study reveals that these methods can perform at random or below under different conditions: semantic entropy achieved AUROC 0.551 on HaluEval but only 0.289 on TruthfulQA in our experiments. This gap between published performance and observed behavior poses a challenge for practitioners selecting detection methods.

Large language models hallucinate confidently, generating plausible-sounding but factually incorrect outputs \citep{lin2022truthfulqa}. Multiple uncertainty-based detection methods have emerged. Semantic entropy clusters generated responses by semantic equivalence via natural language inference (NLI), computing entropy over the cluster distribution \citep{kuhn2023semantic}. Self-consistency measures agreement across multiple generations using surface similarity metrics \citep{manakul2023selfcheckgpt}. Contextual calibration adjusts confidence scores using content-free inputs \citep{zhao2021calibrate}.

Each method has been evaluated on different benchmarks with different sample counts and different models. Semantic entropy reports AUROC on TruthfulQA using specific NLI models; SelfCheckGPT reports on WikiBio using BERTScore. No study, to our knowledge, compares these methods head-to-head on the same benchmarks under matched computational budgets. This leaves practitioners unable to make informed choices---published results are not directly comparable.

This pilot study reveals that hallucination detection method effectiveness may be benchmark-sensitive. Under identical conditions (N=10 samples per query, same model, same temperature), semantic entropy achieved AUROC 0.551 on HaluEval but only 0.289 on TruthfulQA---the latter worse than random, with uncertainty scores inverted relative to ground truth labels. Self-consistency performed near random on both datasets. This suggests that benchmark design may affect method evaluation.

This work makes the following contributions:

\begin{enumerate}
\item \textbf{Matched-Budget Evaluation Framework.} A pilot comparison protocol that tests semantic entropy and self-consistency under identical computational budgets (same sample count, model, and temperature) across multiple benchmarks.
\end{enumerate}

\begin{enumerate}
\item \textbf{Benchmark Sensitivity Observation.} Empirical evidence that detection method performance varies across benchmarks in this pilot---semantic entropy shows marginal signal on HaluEval but inverted results on TruthfulQA.
\end{enumerate}

\begin{enumerate}
\item \textbf{Methodological Observations.} Identification of labeling methodology and benchmark design as potential sources of evaluation inconsistency, providing directions for investigation.
\end{enumerate}
